\documentclass{beamer}
\usepackage[utf8]{inputenc}
\usepackage[russian]{babel}

\title{Настройка балансировщика нагрузки HAProxy}
\subtitle{Курсовой проект по курсу <<Администрирование ОС Linux>>}
\author{Выполнил: студент гр. ХХ-ХХ}
\institute{Кафедра Системного Администрирования}
\date{2026}

\usetheme{Madrid}

\begin{document}

\frame{\titlepage}

\begin{frame}{Цель и задачи проекта}
    \textbf{Цель:} Создание отказоустойчивой и масштабируемой системы веб-серверов с балансировкой нагрузки на базе HAProxy.
    
    \vspace{0.3cm}
    \textbf{Задачи:}
    \begin{itemize}
        \item Настройка HAProxy в режиме Round Robin.
        \item Изоляция инфраструктуры с помощью Docker bridge сетей.
        \item Реализация Sticky Sessions на куках.
        \item Сбор метрик с помощью Prometheus и Grafana.
        \item Написание сценариев автоматизированного Chaos Engineering тестирования.
    end{itemize}
\end{frame}

\begin{frame}{Архитектура решения}
    \begin{block}{Сетевые контуры}
        \begin{itemize}
            \item \textbf{frontend-net:} Только клиент и балансировщик (HAProxy).
            \item \textbf{backend-net:} Изолированный сегмент между HAProxy и тремя бэкендами Nginx.
            \item \textbf{monitoring-net:} Контур сбора метрик (Prometheus, Grafana, cAdvisor).
        \end{itemize}
    \end{block}
\end{frame}

\begin{frame}{Концепция безопасности (Hardening)}
    \begin{itemize}
        \item Запуск всех контейнеров от пользователей без прав суперпользователя (\textbf{non-root}).
        \item Настройка файловой системы бэкендов в режиме \textbf{read-only}.
        \item Монтирование временных папок Nginx во временную файловую систему \textbf{tmpfs}.
        \item Загрузка параметров ядра через \textbf{sysctl} (защита от SYN Flood, IP-spoofing).
    \end{itemize}
\end{frame}

\begin{frame}{Интеграционное тестирование и Chaos Engineering}
    \begin{itemize}
        \item \textbf{Тест Round Robin:} Проверка распределения HTTP запросов.
        \item \textbf{Тест Sticky Sessions:} Проверка удержания сессии клиента по cookie.
        \item \textbf{Сценарий отказа (Chaos):}
        \begin{itemize}
            \item Программный останов 1 или 2 бэкенд серверов.
            \item Замер времени переключения на живой узел.
            \item Проверка отсутствия ошибок HTTP 502/503 на стороне клиента.
        \end{itemize}
    \end{itemize}
\end{frame}

\begin{frame}{Результаты}
    Разработан законченный IaC-репозиторий проекта, позволяющий одной командой развернуть готовую инфраструктуру с балансировкой, мониторингом и системой контроля надежности.
\end{frame}

\end{document}
